\section{Discussion}
\label{sec:discussion}
The experiments support a mechanism-level interpretation of OOB-aware RFFI. OOB evidence is useful because it contains transmitter-discriminative variation, as shown by the signal audit, device probes and the large loss of accuracy when OOB information is replaced or mismatched. The same evidence is risky because OOB-only interventions produce large receiver-level changes while the clean in-band branch and device label remain fixed. The blind replication rules out an explanation based only on one development receiver or one selected seed.

The strong Shen baselines sharpen this interpretation. Since CIS and receiver-adversarial training reach approximately 96--100\% on the blind receivers, receiver shift is not an unavoidable ceiling of the dataset. Instead, different representations expose different deployment behavior. The B1/C$'$ contrast with TrueIB also prevents a universal statement that adding OOB always helps or always hurts. The appropriate conclusion is conditional: the learned dependency of an OOB-aware representation can be receiver-sensitive, architecture-dependent and spectrally asymmetric.

The RSA result defines a practical boundary. A model can become insensitive to one controlled nuisance while retaining poor real receiver transfer. Robustness claims should therefore name the intervention family and should not be promoted to universal receiver invariance without an independent held-out receiver test. In this sense, the blind protocol is part of the scientific result, not merely an evaluation detail.

This study has limitations. The signal-level grouping analysis is descriptive because the released HDF5 files do not provide a complete capture hierarchy. The blind confirmation contains six receiver domains, which is substantial for a fixed public corpus but does not justify universal claims across hardware, modulation or channel families. The interventions probe learned dependencies and do not identify a unique analogue cause. Finally, the task is closed-set device identification, not open-set authentication. Accordingly, we make no claims regarding unknown-device rejection, FAR, FRR, or EER. Future work can test whether the same audit transfers to other waveform families and whether receiver-aware calibration can reduce OOB sensitivity without discarding transmitter information.
